% eksperymentalne tłumaczenie na włoski

%

\subsection{Okrąg opisany i trzy symetralne} % lang-pl

\subsection{Circonferenza circoscritta e tre assi dei segmenti} % lang-it

\begin{proposition}
    \label{symetralne_przecinaja_sie}
    Trzy symetralne boków trójkąta przecinają się w jednym punkcie, środku okręgu opisanego na tym trójkącie. % lang-pl
    
    I tre assi dei lati di un triangolo si intersecano in un unico punto, il centro della circonferenza circoscritta al triangolo. % lang-it
\end{proposition}

Wynika to bezpośrednio z uwagi za definicją \ref{def_symetralna}: w trójkącie $\triangle ABC$, symetralna boku $AB$ (odpowiednio: $AC$) zawiera środki okręgów przechodzących przez punkty $A$ oraz $B$ (przez punkty $A$ oraz $C$), zatem trzecia symetralna musi przejść przez punkt wspólny dwóch pierwszych. % lang-pl
Ciò segue direttamente dall'osservazione successiva alla definizione \ref{def_symetralna}: nel triangolo $\triangle ABC$, l'asse del lato $AB$ (rispettivamente $AC$) contiene i centri delle circonferenze passanti per i punti $A$ e $B$ (rispettivamente per $A$ e $C$), dunque il terzo asse deve passare per il punto comune dei primi due. % lang-it


Trzy różne punkty współliniowe nie leżą na jednym okręgu, zaś dowolne trzy niewspółliniowe punkty leżą na dokładnie jednym okręgu. % lang-pl
 % Guzicki s. 15 % lang-pl
Hartshorne \cite[s. 16]{hartshorne2000} podaje to w formie ćwiczenia ze wskazówką, by spojrzeć na (IV.5). % lang-pl
Audin \cite[s. 61]{audin_2003} też, ale bez wskazówki. % lang-pl
Zetel \cite[s. 151]{zetel_2020} poda ciekawe uzasadnienie z trójkątami rzutowymi. % lang-pl
Inne odsyłacze to Neugebauer \cite[s. 19]{neugebauer_2018}. % lang-pl
Tre punti distinti e allineati non appartengono a una stessa circonferenza, mentre tre punti qualsiasi non allineati appartengono a un'unica circonferenza. % Guzicki s. 15 % lang-it


Hartshorne \cite[p.~16]{hartshorne2000} lo propone come esercizio, suggerendo di consultare (IV.5). % lang-it
Anche Audin \cite[p.~61]{audin_2003}, ma senza suggerimenti. % lang-it
Zetel \cite[p.~151]{zetel_2020} ne fornisce una dimostrazione interessante mediante triangoli proiettivi. % lang-it
Un altro riferimento è Neugebauer \cite[p.~19]{neugebauer_2018}. % lang-it
\begin{proposition}
	Środek okręgu opisanego na trójkącie leży wewnątrz tego trójkąta (odpowiednio: na jego brzegu, na zewnątrz trójkąta) wtedy i tylko wtedy, gdy trójkąt jest ostrokątny (odpowiednio: prostokątny, rozwartokątny). % lang-pl
	
	Il centro della circonferenza circoscritta a un triangolo si trova all'interno del triangolo (rispettivamente sul suo bordo, all'esterno del triangolo) se e solo se il triangolo è acutangolo (rispettivamente rettangolo, ottusangolo). % lang-it
\end{proposition}

Uogólnieniem faktu \ref{symetralne_przecinaja_sie} o współpękowości symetralnych boków trójkąta jest: % lang-pl

Una generalizzazione del fatto \ref{symetralne_przecinaja_sie} sulla concorrenza degli assi dei lati di un triangolo è il seguente risultato: % lang-it

\begin{theorem}[Carnota]
\index{twierdzenie!Carnota}%
\index{teorema!di Carnot}% % lang-it
\label{guzicki_6_13}%
	Dany jest trójkąt $ABC$ i punkty $D, E, F$ leżące odpowiednio na prostych $BC, CA, AB$. % lang-pl
	
	Niech prosta $k$ (odpowiednio: $l$, $m$) przechodzi przez punkt $D$ ($E$, $F$) i będzie prostopadła do prostej $BC$ ($CA$, $AB$). % lang-pl
	Siano dati un triangolo $ABC$ e punti $D$, $E$, $F$ appartenenti rispettivamente alle rette $BC$, $CA$, $AB$. % lang-it
	
	Wtedy proste $k$, $l$, $m$ mają punkt wspólny wtedy i tylko wtedy, gdy % lang-pl
	Sia la retta $k$ (rispettivamente $l$, $m$) passante per $D$ ($E$, $F$) e perpendicolare alla retta $BC$ ($CA$, $AB$). % lang-it
	
	Allora le rette $k$, $l$, $m$ sono concorrenti se e solo se % lang-it
	\begin{equation}
		|AF|^2 + |BD|^2 + |CE|^2 = |AE|^2 + |BF|^2 + |CD|^2.
	\end{equation}
\end{theorem}
% TODO: https://en.wikipedia.org/wiki/Carnot%27s_theorem_(perpendiculars)

Guzicki \cite[s. 176]{guzicki_2021} wyprowadza je z twierdzenia Pitagorasa, co pozwala mu dojść do wniosku, że okręgi opisany i wpisany oraz ortocentrum istnieją. % lang-pl

Guzicki \cite[p.~176]{guzicki_2021} lo deduce dal teorema di Pitagora, ottenendo così l'esistenza della circonferenza circoscritta, della circonferenza inscritta e dell'ortocentro. % lang-it
\index{twierdzenie!Pitagorasa}
\index{teorema!di Pitagora} % lang-it

%